% First draft — embedding-based dependence testing between graph populations.
% Build: pdflatex main && bibtex main && pdflatex main && pdflatex main
\documentclass[11pt]{article}

\usepackage[margin=1.1in]{geometry}
\usepackage{amsmath,amssymb,amsthm}
\usepackage{booktabs}
\usepackage{graphicx}
\usepackage[colorlinks=true,linkcolor=blue,citecolor=blue,urlcolor=blue]{hyperref}
\usepackage[round]{natbib}
\usepackage{xcolor}

\newcommand{\todo}[1]{\textcolor{red}{[TODO: #1]}}
\newcommand{\dCor}{\mathrm{dCor}}
\newtheorem{remark}{Remark}
\newtheorem{proposition}{Proposition}

\title{Graph Embeddings as Summary Statistics for Testing\\
Dependence Between Graph Populations}

\author{Maruan Ottoni Bakri\thanks{maruanbakriottoni@gmail.com. Code and full
experimental logs: \url{https://github.com/mbottoni/graph-vae-subsampling}.}}

\date{June 2026 --- working draft}

\begin{document}
\maketitle

\begin{abstract}
Measuring statistical dependence between sets of graphs is a central problem in
domains such as brain-network analysis, where the main difficulty is mapping
graph structure into vectors amenable to classical tests. The state-of-the-art
approach correlates the spectral radii of paired graphs
\citep{fujita2017correlation}. We show that this test --- and every \emph{raw}
spectral statistic we evaluate, including the full eigenvalue spectrum, NetLSD
heat traces, and kernel (HSIC) variants --- fails in a simple and practically
relevant regime: stochastic block model populations whose dependence is
confined to the between-block connection probability $p_{\mathrm{out}}$
(power $\le 0.19$ at $\alpha = 0.05$, $R = 100$). The failure is
\emph{geometric}, not informational: the spectrum contains the dependence
signal in closed form ($(\lambda_1 - \lambda_2)/2$ tracks $p_{\mathrm{out}}$ at
$r = 0.89$), but distance-based dependence measures are dominated by
independent nuisance variation in $p_{\mathrm{in}}$, which every leading
eigenvalue loads on. Two statistics restore power to $\approx 0.9$: (i) a
\emph{whitened} truncated spectrum ($0.97$) --- the sample analogue of
affinely invariant distance correlation --- and (ii) the singular values of
the latent node-embedding matrix of a per-graph variational graph autoencoder
(VGAE) ($0.89$), whose leading coordinate aligns with $p_{\mathrm{out}}$
\emph{without} explicit nuisance normalization. The gain is not free: the same
whitening costs power when dependence is spectrum-aligned, and a Bonferroni
combiner is the most robust single choice across regimes. All tests are
exactly calibrated under permutation nulls. The headline message is that
representation geometry, not information content, decides test power for graph
dependence.
\todo{rerun at $R{=}100$; real fMRI data; combined statistic}
\end{abstract}

\section{Introduction}

Graphs are rich representations used across biology, neuroscience, and the
social sciences. A recurring statistical task is to decide whether two
\emph{populations} of graphs are dependent: given paired observations
$\{(G^{(1)}_i, G^{(2)}_i)\}_{i=1}^{k}$ --- for instance, two brain networks
recorded per subject --- test
\[
H_0: G^{(1)} \perp G^{(2)}
\quad\text{vs.}\quad
H_1: G^{(1)} \not\perp G^{(2)},
\]
where dependence is induced at the level of the generative process rather than
through any vertex correspondence. The central obstacle is the translation of
graph components (e.g.\ the adjacency matrix) into vectors of values on which a
dependence test can operate.

\citet{fujita2017correlation} proposed the state-of-the-art solution: map each
graph to its spectral radius $\lambda_{\max}(A)$ --- a quantity that concentrates
sharply around a function of the generative parameters for many random-graph
models --- and test for correlation between the resulting paired scalar series.
The spectral radius summarizes important structural properties such as cliques
and walk counts, and the resulting test is simple, calibrated, and effective
when the dependence signal is aligned with it.

Meanwhile, the graph-embedding literature has produced flexible learned
mappings from graphs into latent vector spaces that preserve structural
properties \citep{cai2018survey, kipf2016variational}, optionally incorporating
node and edge features. Despite their success in prediction tasks, learned
embeddings have not been explored as \emph{summary statistics for dependence
testing between graph populations}. This work fills that gap.

\paragraph{Contributions.}
\begin{enumerate}
  \item We formulate embedding-based dependence testing for graph populations:
        map each graph to a rotation-invariant summary of a learned embedding,
        then apply a permutation distance-correlation test
        (Section~\ref{sec:method}).
  \item We show that all statistics considered --- the spectral-radius benchmark,
        truncated spectra, structural features, and learned-embedding summaries
        --- are exactly calibrated under permutation nulls
        (Section~\ref{sec:calibration}).
  \item On stochastic block model (SBM) populations whose dependence is
        confined to the between-block density $p_{\mathrm{out}}$, the benchmark
        and every \emph{raw} spectral statistic fail (power $\le 0.19$ at
        $R = 100$, $\le 0.36$ at $R = 25$), including the full spectrum,
        NetLSD, and HSIC variants (Section~\ref{sec:sbm}).
  \item We identify the mechanism and verify it twice. Adjacency eigenvalues
        \emph{entangle} the SBM parameters as $\lambda_{1,2} \approx
        \tfrac{n}{2}(p_{\mathrm{in}} \pm p_{\mathrm{out}})$, so distance-based
        dependence measures on spectra are dominated by independent nuisance
        variation in $p_{\mathrm{in}}$ --- \emph{even though}
        $(\lambda_1 - \lambda_2)/2$ isolates the signal exactly (measured:
        $r = 0.89$ with $p_{\mathrm{out}}$, $r = 0.15$ with $p_{\mathrm{in}}$).
        Confirmations: (i) \emph{whitening} the truncated spectrum --- a purely
        geometric correction --- restores power from $0.19$ to $0.97$; (ii) the
        VGAE latent geometry disentangles the parameters (its leading singular
        value correlates $-0.61$ with $p_{\mathrm{out}}$ but only $0.19$ with
        $p_{\mathrm{in}}$) and reaches $0.89$ without any explicit correction
        (Section~\ref{sec:mechanism}).
  \item Practical guidance: there is no uniformly best statistic --- whitening
        trades power in the aligned regime for power in the blind spot --- but a
        Bonferroni $\min$-$p$ combiner of the whitened spectrum and the learned
        summary is the most robust single choice across regimes
        (Section~\ref{sec:sbm}).
\end{enumerate}

\section{Related Work}\label{sec:related}

\paragraph{Dependence between graphs.} \citet{fujita2017correlation}
introduced the spectral-radius correlation test for populations of paired
graphs; it is our benchmark and defines our problem setting. A distinct line of
work tests independence between two \emph{vertex-aligned} graphs observed
once: \citet{xiong2019graphindep} study a single pair of $n \times n$
adjacency matrices under a $\rho$-correlated SBM, show that na\"ive
permutation tests are invalid there, and propose block-permutation nulls. Our
setting differs on both axes --- a population of $k$ paired graphs with no
vertex alignment, dependence induced at the generative-parameter level ---
which is why plain permutation of the pairing is exact here. Graph two-sample
testing and graph-kernel MMD methods \citep{gretton2012kernel} compare
distributions of graphs rather than testing dependence between paired
populations, but supply natural baseline summaries.

\paragraph{Graph embeddings.} Survey: \citet{cai2018survey}. We build on the
variational graph autoencoder \citep{kipf2016variational}. Spectral
descriptors such as NetLSD \citep{tsitsulin2018netlsd} are fixed (non-learned)
alternatives included in our baseline set.

\paragraph{Affinely invariant dependence measures.} Whitening before dCor is
the sample analogue of affinely invariant distance correlation
\citep{dueck2014affinely}. The whitening idea is therefore not ours; the
contributions are the diagnosis that the graph-population benchmark fails for
geometric (not informational) reasons, the demonstration on graph spectra, and
the observation that learned embeddings achieve the same effect without
explicit normalization.

\section{Method}\label{sec:method}

\subsection{Setting}

For pair $i$, latent generative parameters $\theta_i^{(1)}, \theta_i^{(2)}$ are
drawn from a joint distribution with marginals fixed and dependence controlled
by $\rho \in [0, 1]$; graphs are then drawn independently given the parameters:
$G_i^{(j)} \sim \mathbb{P}(\cdot \mid \theta_i^{(j)})$. Dependence lives
entirely at the parameter level, which is what a summary statistic must expose.
We use the \emph{shared-parameter construction}: $\theta_i^{(2)} =
\theta_i^{(1)}$ with probability $\rho$, otherwise an independent draw. This
keeps marginals exact and gives $\mathrm{corr}(\theta^{(1)}, \theta^{(2)}) =
\rho$.

\subsection{Summary statistics}

Each graph $G$ is mapped to a vector $s(G) \in \mathbb{R}^q$:

\begin{description}
  \item[\texttt{lambda\_max} (benchmark).] $s(G) = \lambda_{\max}(A)$
        \citep{fujita2017correlation}.
  \item[\texttt{spectral}$q$.] Top-$q$ adjacency eigenvalues, descending
        ($q = 5$ throughout).
  \item[\texttt{features}.] Density, mean clustering, normalized degree mean
        and standard deviation, and extreme eigenvalues.
  \item[\texttt{vgae\_sv} (proposed).] Train a VGAE \citep{kipf2016variational}
        on $G$ alone (two-layer GCN encoder, inner-product decoder, full-graph
        reconstruction objective) to obtain the latent matrix $Z \in
        \mathbb{R}^{n \times d}$; return the top-$q$ singular values of the
        column-centered $Z$, scaled by $n^{-1/2}$.
\end{description}

\begin{remark}[Rotation invariance]\label{rem:rotation}
Per-graph training makes the latent axes arbitrary: any orthogonal rotation of
$Z$ achieves the same inner-product reconstruction. Raw latent coordinates are
therefore not comparable across graphs. The singular values of $Z$ are
invariant to this rotation and constitute the natural analogue of the
adjacency spectrum --- making \texttt{vgae\_sv} vs.\ \texttt{spectral}$q$ a
clean comparison between a \emph{learned} and a \emph{fixed} spectrum.
\end{remark}

\subsection{Test}

Scalar summaries are tested with $|$Pearson $r$$|$; vector summaries with
distance correlation \citep{szekely2007measuring}, which is invariant to
orthogonal transformations of either space. All statistics use permutation
null distributions ($B = 200$ permutations of the pairing), so every test is
exact at the same level $\alpha = 0.05$ and power comparisons are like-for-like.

\section{Experiments}\label{sec:experiments}

Throughout: $k = 40$ pairs per replicate, $n = 60$ nodes per graph, $R$
replicates per configuration, rejection rates at $\alpha = 0.05$. VGAE: hidden
width 16, latent dimension 8, 60 epochs, Adam ($\mathrm{lr} = 0.01$).
\todo{replace $R \le 25$ tables with the definitive $R = 100$ run (in flight);
sensitivity in $n \in \{60, 100, 200\}$, $k \in \{20, 40, 80\}$}

\paragraph{A cautionary note on simulation harnesses.}
An early version of our harness derived all randomness in a replicate --- the
data-generating RNG, the per-graph VGAE training seeds, and the permutation
RNG --- from a single base seed, with consecutive replicates using consecutive
seeds. This coupling produced an apparent type-I rate of $0.10$
\mbox{[$0.06, 0.17$]} for the learned statistic at $R = 100$, indistinguishable
from a defect in the statistic itself. Re-running the pure-null cell with three
mutually independent \texttt{SeedSequence}-spawned streams restored exactness:
type-I $0.050$ for the learned statistic ($0.030$ for the benchmark and the
whitened spectrum, $0.020$ for the Bonferroni combination --- conservative, as
expected). All headline tables use the decoupled harness. Permutation tests
are exact only when the permutation randomness is independent of the data;
simulation code can silently violate this assumption at the harness level
while every individual component is correct.

\subsection{Calibration and the one-parameter regime (ER)}\label{sec:calibration}

Pairs of Erd\H{o}s--R\'enyi graphs with $p_i \sim U(0.1, 0.3)$ under the
shared-parameter construction ($R = 20$).

\begin{table}[ht]
\centering
\caption{Rejection rates, ER pairs. The $\rho = 0$ row is type-I error.}
\label{tab:er}
\begin{tabular}{lcccc}
\toprule
$\rho$ & \texttt{lambda\_max} & \texttt{spectral5} & \texttt{features} & \texttt{vgae\_sv} \\
\midrule
0.00 & 0.05 & 0.05 & 0.05 & 0.05 \\
0.25 & 0.25 & 0.25 & 0.20 & 0.35 \\
0.50 & 0.90 & 0.95 & 0.95 & 0.80 \\
0.75 & 1.00 & 1.00 & 1.00 & 1.00 \\
1.00 & 1.00 & 1.00 & 1.00 & 1.00 \\
\bottomrule
\end{tabular}
\end{table}

All four tests are exactly calibrated (rejection $0.05$ at $\rho = 0$). Power
is comparable across statistics: ER is the benchmark-friendliest setting, since
a single parameter drives all structure and $\lambda_{\max} \approx np$ is
nearly sufficient for it. Differences at low $\rho$ are within Monte-Carlo
noise at $R = 20$.

\subsection{Two-parameter dependence: the SBM blind spot}\label{sec:sbm}

Pairs of two-block SBMs ($n = 60$, blocks of $30$) with
$p_{\mathrm{in}, i} \sim U(0.15, 0.35)$ and
$p_{\mathrm{out}, i} \sim U(0.02, 0.10)$. Two conditions:
\textbf{both} --- each parameter shared with probability $\rho$;
\textbf{$p_{\mathrm{out}}$-only} --- only $p_{\mathrm{out}}$ correlated,
$p_{\mathrm{in}}$ independent across the pair. Table~\ref{tab:sbm} reports the
definitive run: $R = 100$ replicates per cell with the decoupled seeding of
Section~\ref{sec:experiments}, $95\%$ Wilson intervals in brackets. We include
two combiners of the whitened spectrum and the learned summary: their
concatenation (\texttt{concat}) and the Bonferroni $\min$-$p$ test
(\texttt{bonf}, reject if $2\min(p_{\mathrm{sw}}, p_{\mathrm{vgae}}) < \alpha$).

\begin{table}[ht]
\centering
\small
\caption{Rejection rates, SBM pairs, $R = 100$ ($95\%$ Wilson CI). Top:
dependence in both parameters (spectrum-aligned). Bottom: dependence only in
$p_{\mathrm{out}}$ (the blind spot). The $\rho = 0$ rows are type-I error: all
calibrated. \texttt{sw} = whitened top-5 spectrum, \texttt{vg} = VGAE singular
values, \texttt{bonf} = Bonferroni $\min$-$p$.}
\label{tab:sbm}
\begin{tabular}{llcccccc}
\toprule
cond. & $\rho$ & $\lambda_{\max}$ & spec5 & sw & vg & concat & bonf \\
\midrule
both & 0.0 & .00 & .01 & .04 & .05 & .06 & .04 \\
both & 0.5 & \textbf{.90} & .87 & .65 & .35 & .64 & .60 \\
both & 1.0 & 1.0 & 1.0 & 1.0 & .92 & 1.0 & 1.0 \\
\addlinespace
$p_{\mathrm{out}}$ & 0.0 & .07 & .04 & .04 & .07 & .07 & .05 \\
$p_{\mathrm{out}}$ & 0.5 & .06 & .07 & .18 & .30 & .21 & \textbf{.33} \\
$p_{\mathrm{out}}$ & 1.0 & .07 & .19 & .97 & .89 & \textbf{.99} & .98 \\
\bottomrule
\end{tabular}
\end{table}

Three facts. (i) \emph{The benchmark has a blind spot:} in the
$p_{\mathrm{out}}$-only condition the spectral radius never exceeds $0.07$ even
at $\rho = 1$, and the raw top-5 spectrum reaches only $0.19$
\mbox{[$0.13, 0.28$]} --- both indistinguishable from noise. (ii) \emph{Two
statistics see it:} the whitened spectrum ($0.97$) and the learned summary
($0.89$) both detect the dependence, with the whitened spectrum slightly ahead
(overlapping CIs at $\rho = 1$, separated at $\rho = 0.5$). (iii) \emph{There
is no free lunch:} the same whitening that rescues the blind spot \emph{costs}
power when dependence is spectrum-aligned ($0.65$ vs.\ the raw $0.90$ at
$\rho = 0.5$, \textbf{both}), and the learned summary costs more ($0.35$). The
Bonferroni combiner is the most robust single choice --- best of all methods in
the hardest cell ($0.33$ at $p_{\mathrm{out}}$, $\rho = 0.5$, beating both of
its components) and near-top in the blind spot ($0.98$) --- at a moderate price
in the aligned regime ($0.60$). All type-I rates ($\rho = 0$) lie in
$[0.00, 0.07]$.

\subsection{Stronger baselines: more information does not help; better geometry does}
\label{sec:killorconfirm}

Could the failure in Table~\ref{tab:sbm} be an artifact of truncating the
spectrum at five eigenvalues, of using dCor rather than a kernel statistic, or
of ignoring multi-scale structure? Table~\ref{tab:d2b} ($R = 25$) adds the
baselines that rule these out: dCor on \emph{all} $n = 60$ eigenvalues
(\texttt{full\_spec}), NetLSD heat-trace signatures \citep{tsitsulin2018netlsd}
(\texttt{netlsd}), HSIC with Gaussian kernels and median-heuristic bandwidth
(\texttt{hsic5}), and dCor on the ZCA-\emph{whitened} top-5 spectrum
(\texttt{spec5w}) --- the sample analogue of affinely invariant distance
correlation \citep{dueck2014affinely}, and the direct test of the geometric
explanation: whitening changes no information, only the metric.

\begin{table}[ht]
\centering
\caption{Kill-or-confirm baselines, SBM pairs ($R = 25$, identical seeds and
protocol across columns). Adding information (\texttt{full\_spec},
\texttt{netlsd}) or changing the dependence measure (\texttt{hsic5}) does not
repair the $p_{\mathrm{out}}$-only failure; fixing the geometry
(\texttt{spec5w}) does, matching the learned summary. The $R = 100$ figures for
the overlapping columns appear in Table~\ref{tab:sbm}.}
\label{tab:d2b}
\begin{tabular}{llccccccc}
\toprule
cond. & $\rho$ & $\lambda_{\max}$ & spec5 & full\_spec & netlsd & hsic5 & spec5w & vgae \\
\midrule
both        & 0.0 & 0.00 & 0.00 & 0.00 & 0.00 & 0.00 & 0.08 & 0.04 \\
both        & 0.5 & 0.84 & 0.88 & 0.84 & 0.80 & 0.84 & 0.76 & 0.24 \\
both        & 1.0 & 1.00 & 1.00 & 1.00 & 1.00 & 1.00 & 1.00 & 0.92 \\
\addlinespace
$p_{\mathrm{out}}$ & 0.0 & 0.00 & 0.00 & 0.00 & 0.00 & 0.00 & 0.00 & 0.08 \\
$p_{\mathrm{out}}$ & 0.5 & 0.08 & 0.16 & 0.16 & 0.12 & 0.12 & \textbf{0.28} & 0.16 \\
$p_{\mathrm{out}}$ & 1.0 & 0.24 & 0.24 & 0.36 & 0.32 & 0.28 & \textbf{0.96} & \textbf{0.96} \\
\bottomrule
\end{tabular}
\end{table}

Three conclusions. First, the failure is not informational: sixty eigenvalues
do no better than one ($0.36$ vs.\ $0.24$), and neither a multi-scale
descriptor nor a kernel statistic helps. Second, the failure is geometric: a
linear re-metrization of the same five numbers lifts power from $0.24$ to
$0.96$. Third, the learned summary reaches the same power \emph{without}
explicit whitening --- but pays for it when dependence is aligned with the
dominant spectral directions ($0.24$ vs.\ $0.76$--$0.88$ at $\rho = 0.5$,
\textbf{both}). The whitened spectrum is therefore the practical
recommendation; the learned summary is the demonstration that representation
geometry is the deciding quantity.

\subsection{Sensitivity to graph and population size}\label{sec:sensitivity}

Does the blind spot, and its repair, survive scaling? Table~\ref{tab:sens}
sweeps graph size $n$ and population size $k$ in the $p_{\mathrm{out}}$-only
condition at $\rho = 1$ ($R = 100$, fixed before inspection by three
pre-registered predictions). All three hold: (P1) the benchmark remains blind
for every $n$ ($\lambda_{\max}$ power $0.11$--$0.17$ at $k = 40$, flat); (P2)
the whitening rescue is monotone in $n$, saturating to $1.00$ by $n = 100$ at
$k = 40$, consistent with the eigenvalue concentration $O(\sqrt{n})$ vs.\ signal
$O(n)$; (P3) the benchmark-to-whitening gap widens with $k$ until whitening
reaches the ceiling.

\begin{table}[ht]
\centering
\small
\caption{Power in the $p_{\mathrm{out}}$ blind spot at $\rho = 1$, $R = 100$,
across $n$ (rows) and $k$ (cols). The benchmark ($\lambda_{\max}$) is blind
everywhere; the whitened spectrum (\texttt{sw}) is rescued, faster as $n$ and
$k$ grow. The raw truncated spectrum (\texttt{s5}) is blind only at small
samples --- its higher eigenvalues expose $p_{\mathrm{out}}$ once concentration
and sample size suffice ($0.58$ at $n = 150, k = 80$) --- showing the
benchmark's blindness is \emph{structural} (one number cannot isolate
$p_{\mathrm{out}}$) while the spectrum's is merely finite-sample, and whitening
buys the same signal far sooner.}
\label{tab:sens}
\begin{tabular}{ll ccc}
\toprule
 & & $k=20$ & $k=40$ & $k=80$ \\
\midrule
$n=40$  & $\lambda_{\max}$ / s5 / sw & .07/.09/.18 & .14/.12/.61 & .27/.18/.98 \\
$n=60$  & $\lambda_{\max}$ / s5 / sw & .11/.16/.38 & .11/.12/.94 & .30/.36/1.0 \\
$n=100$ & $\lambda_{\max}$ / s5 / sw & .16/.12/.66 & .12/.14/1.0 & .25/.46/1.0 \\
$n=150$ & $\lambda_{\max}$ / s5 / sw & .06/.05/.69 & .17/.33/1.0 & .20/.58/1.0 \\
\bottomrule
\end{tabular}
\end{table}

The finding refines the headline: the spectral \emph{radius} has a structural
blind spot no sample size cures, whereas the full truncated spectrum is only
underpowered --- but whitening reaches saturation while the raw spectrum is
still at chance (e.g.\ $n = 60, k = 40$: $0.94$ vs.\ $0.12$). (At $k = 20$ the
benchmark's permutation null is itself unstable, type-I up to $0.13$; $k \ge 40$
is calibrated.)

\subsection{Mechanism: entanglement vs.\ disentanglement}\label{sec:mechanism}

For a balanced two-block SBM the leading adjacency eigenvalues concentrate
around
\[
\lambda_1 \approx \tfrac{n}{2}\left(p_{\mathrm{in}} + p_{\mathrm{out}}\right),
\qquad
\lambda_2 \approx \tfrac{n}{2}\left(p_{\mathrm{in}} - p_{\mathrm{out}}\right).
\]
The dependence signal in the $p_{\mathrm{out}}$-only condition is therefore
\emph{available in closed form} from the spectrum, as
$(\lambda_1 - \lambda_2)/2 \approx \tfrac{n}{2}\, p_{\mathrm{out}}$ --- yet
\texttt{spectral5} fails ($0.24$). The explanation is geometric: distance
correlation operates on inter-point distances, and in the spectral
representation those distances are dominated by the independent
$p_{\mathrm{in}}$ variation (range $0.20$) relative to the correlated
$p_{\mathrm{out}}$ variation (range $0.08$), which enters both leading
eigenvalues with equal magnitude. Information content is not sufficient for
test power; the \emph{orientation} of the summary's coordinates relative to
nuisance variation decides it.

We verify the entanglement claim directly: over $200$ independent SBMs sampled
from the experimental parameter ranges, every one of the five leading
adjacency eigenvalues correlates dominantly with $p_{\mathrm{in}}$
($r = 0.63$--$0.92$), with $\lambda_1$ and $\lambda_2$ loading on
$p_{\mathrm{out}}$ with \emph{opposite signs} ($+0.32$ / $-0.31$), exactly as
the $\tfrac{n}{2}(p_{\mathrm{in}} \pm p_{\mathrm{out}})$ approximation
predicts. The derived combination $(\lambda_1 - \lambda_2)/2$ isolates the
signal ($r = 0.89$ with $p_{\mathrm{out}}$, $r = 0.15$ with $p_{\mathrm{in}}$)
--- present in the spectrum, aligned with no raw coordinate
(Table~\ref{tab:mech}).

\begin{table}[ht]
\centering
\caption{Pearson correlation of each summary coordinate with the generative
parameters ($200$ independent SBMs). Eigenvalues entangle the parameters; the
leading learned singular value is the $p_{\mathrm{out}}$ coordinate.}
\label{tab:mech}
\begin{tabular}{lcc@{\qquad}lcc}
\toprule
coordinate & $p_{\mathrm{in}}$ & $p_{\mathrm{out}}$ & coordinate & $p_{\mathrm{in}}$ & $p_{\mathrm{out}}$ \\
\midrule
$\lambda_1$ & 0.89 & 0.32  & $\mathrm{sv}_1$ & 0.19 & $-$0.61 \\
$\lambda_2$ & 0.92 & $-$0.31 & $\mathrm{sv}_2$ & $-$0.59 & $-$0.30 \\
$\lambda_3$ & 0.67 & 0.43  & $\mathrm{sv}_3$ & $-$0.75 & $-$0.23 \\
$\lambda_4$ & 0.68 & 0.46  & $\mathrm{sv}_4$ & $-$0.73 & $-$0.16 \\
$\lambda_5$ & 0.63 & 0.55  & $\mathrm{sv}_5$ & $-$0.73 & $-$0.11 \\
\midrule
$(\lambda_1{-}\lambda_2)/2$ & 0.15 & 0.89 & & & \\
\bottomrule
\end{tabular}
\end{table}

The VGAE latent geometry, by contrast, disentangles the parameters. The
inner-product decoder places same-block nodes near one another and the two
blocks at an angle determined by the cross-block connection probability:
between-cluster separation tracks $p_{\mathrm{out}}$ while within-cluster
spread tracks $p_{\mathrm{in}}$. Empirically (Table~\ref{tab:mech}), the
\emph{leading} singular value is the $p_{\mathrm{out}}$ coordinate ($r =
-0.61$ vs.\ $0.19$ with $p_{\mathrm{in}}$), with $p_{\mathrm{in}}$ relegated
to $\mathrm{sv}_3$--$\mathrm{sv}_5$. Distance-based dependence measures weight
directions by variance; in the learned representation the highest-variance
direction \emph{is} the signal direction, while in the raw spectrum it is the
nuisance. This explains all of Table~\ref{tab:d2b}: whitening equalizes the
directions (rescuing the fixed spectrum), and the learned summary arrives with
the right orientation built in.

The following idealization makes the orientation claim precise.

\begin{proposition}[Two-centroid latents]\label{prop:centroids}
Consider a balanced two-block SBM with parameters $(p_{\mathrm{in}},
p_{\mathrm{out}})$ and an idealized latent configuration in which block-$1$
nodes sit at $+\mu \in \mathbb{R}^d$ and block-$2$ nodes at $-\mu$, plus
isotropic within-cluster noise of scale $\sigma$, decoded by
$\Pr(A_{ij} = 1) = \operatorname{sigmoid}(z_i^{\top} z_j + c)$ with a shared
offset $c$. Matching expected within- and between-block densities gives
\[
\|\mu\|^2 \;=\; \tfrac{1}{2}\bigl(\operatorname{logit} p_{\mathrm{in}} -
\operatorname{logit} p_{\mathrm{out}}\bigr),
\qquad
c \;=\; \tfrac{1}{2}\bigl(\operatorname{logit} p_{\mathrm{in}} +
\operatorname{logit} p_{\mathrm{out}}\bigr),
\]
so the leading singular value of the centered latent matrix is
$\mathrm{sv}_1 = \sqrt{n}\,\|\mu\| =
\sqrt{\tfrac{n}{2}\bigl(\operatorname{logit} p_{\mathrm{in}} -
\operatorname{logit} p_{\mathrm{out}}\bigr)}$, while the remaining singular
values are $O(\sigma\sqrt{n})$, carrying no direct dependence on the logit
gap. Since $\partial \operatorname{logit}(p) / \partial p = 1/\bigl(p(1 -
p)\bigr)$ is $\approx 3.3\times$ larger at the midpoint of the experimental
$p_{\mathrm{out}}$ range ($0.06$) than of the $p_{\mathrm{in}}$ range
($0.25$), $\mathrm{sv}_1$ is dominated by $p_{\mathrm{out}}$.
\end{proposition}

The idealization is consistent with measurement: over $120$ independent SBMs,
$\mathrm{sv}_1$ correlates $0.645$ with the logit gap ($R^2 = 0.42$), and
after regressing the gap out, the residual correlations with
$p_{\mathrm{in}}$ and $p_{\mathrm{out}}$ collapse to $-0.18$ and $-0.10$ ---
i.e.\ essentially all of $\mathrm{sv}_1$'s parameter dependence flows through
the predicted quantity. (The trailing singular values' empirical
$p_{\mathrm{in}}$ loading, Table~\ref{tab:mech}, corresponds to the noise
scale $\sigma$ tracking within-block density --- outside the idealization but
harmless to it: those coordinates carry nuisance, which dCor down-weights only
if, as here, their variance is smaller.) \todo{tighten the proposition's
noise model; current form is a matched-moments idealization, not a statement
about the VGAE optimum}

\section{Limitations and Next Steps}\label{sec:limitations}

\textbf{Scale.} The headline (Table~\ref{tab:sbm}) is at $R = 100$ with
calibrated type-I, and the blind spot and its repair are confirmed across
$n \in \{40, \dots, 150\}$ and $k \in \{20, \dots, 80\}$
(Section~\ref{sec:sensitivity}). Larger graphs ($n \ge 500$) and the
learned summary's behavior across that grid (only spot-checked here) remain
open.
\textbf{Whitening stability.} ZCA whitening with $k = 40$ samples in $5$
dimensions is benign; in higher-dimensional summaries the covariance estimate
degrades and the learned summary's built-in orientation may become the more
robust option --- untested.
\textbf{Real data.} A paired brain-network application (the benchmark's home
domain) is required for the full claim.
\textbf{Features.} The original program includes testing whether node features
add power (phase two); untested here.
\textbf{Theory.} The disentanglement mechanism is currently an empirical
observation with a heuristic argument.

\section*{Reproducibility}

All experiments are scripted (\texttt{experiments/}) with fixed seeds and JSON
result logs (\texttt{results/}) in the repository; the package
(\texttt{src/gvs}) implements generators, summaries, and permutation tests.

\bibliographystyle{plainnat}
\begin{thebibliography}{9}

\bibitem[Cai et~al.(2018)]{cai2018survey}
H.~Cai, V.~W. Zheng, and K.~C.-C. Chang.
\newblock A comprehensive survey of graph embedding: Problems, techniques, and
  applications.
\newblock \emph{IEEE Transactions on Knowledge and Data Engineering},
  30(9):1616--1637, 2018.

\bibitem[Dueck et~al.(2014)]{dueck2014affinely}
J.~Dueck, D.~Edelmann, T.~Gneiting, and D.~Richards.
\newblock The affinely invariant distance correlation.
\newblock \emph{Bernoulli}, 20(4):2305--2330, 2014.

\bibitem[Fujita et~al.(2017)]{fujita2017correlation}
A.~Fujita, D.~Y. Takahashi, J.~R. Balardin, M.~C. Vidal, and J.~R. Sato.
\newblock Correlation between graphs with an application to brain network
  analysis.
\newblock \emph{Computational Statistics \& Data Analysis}, 109:76--92, 2017.

\bibitem[Gretton et~al.(2012)]{gretton2012kernel}
A.~Gretton, K.~M. Borgwardt, M.~J. Rasch, B.~Sch\"olkopf, and A.~Smola.
\newblock A kernel two-sample test.
\newblock \emph{Journal of Machine Learning Research}, 13:723--773, 2012.

\bibitem[Kipf and Welling(2016)]{kipf2016variational}
T.~N. Kipf and M.~Welling.
\newblock Variational graph auto-encoders.
\newblock \emph{NIPS Workshop on Bayesian Deep Learning}, 2016.

\bibitem[Sz\'ekely et~al.(2007)]{szekely2007measuring}
G.~J. Sz\'ekely, M.~L. Rizzo, and N.~K. Bakirov.
\newblock Measuring and testing dependence by correlation of distances.
\newblock \emph{The Annals of Statistics}, 35(6):2769--2794, 2007.

\bibitem[Tsitsulin et~al.(2018)]{tsitsulin2018netlsd}
A.~Tsitsulin, D.~Mottin, P.~Karras, A.~Bronstein, and E.~M\"uller.
\newblock NetLSD: Hearing the shape of a graph.
\newblock \emph{KDD}, 2018.

\bibitem[Xiong et~al.(2019)]{xiong2019graphindep}
J.~Xiong, C.~Shen, J.~Arroyo, and J.~T. Vogelstein.
\newblock Community correlations and testing independence between binary
  graphs.
\newblock \emph{arXiv:1906.03661}, 2019.

\end{thebibliography}

\end{document}
